\chapter{Explicabilité et rapport de recommandations}\label{chapExpl}
\pagestyle{fancy}
\lhead{Explicabilité et rapport de recommandations}

\graphicspath{{Chapitre4/images_chapitre4/}{images_chapitre4/}{Chapitre3/images_chapitre3/}{images_chapitre3/}}

% Couleurs de tableaux (déjà définies dans le préambule principal : ne sont créées que si absentes)
\providecolor{thead}{HTML}{1F3A5F}
\providecolor{trowa}{HTML}{EEF2F7}

% Insère une figure si le fichier existe, sinon un cadre « à insérer »
\providecommand{\figtrois}[2]{%
  \IfFileExists{Chapitre3/images_chapitre3/#1}{\includegraphics[width=#2\textwidth]{#1}}{%
  \IfFileExists{images_chapitre3/#1}{\includegraphics[width=#2\textwidth]{#1}}{%
  \IfFileExists{Chapitre4/images_chapitre4/#1}{\includegraphics[width=#2\textwidth]{#1}}{%
  \IfFileExists{images_chapitre4/#1}{\includegraphics[width=#2\textwidth]{#1}}{%
  \fbox{\parbox[c][4cm][c]{#2\textwidth}{\centering\small\itshape
  Figure à insérer : \texttt{\detokenize{#1}}}}}}}}}

% =====================================================================
\section{Introduction}

Un modèle qui répond seulement « attaque » ou « normal » ne suffit pas à un
opérateur de chaîne du froid : avant de bloquer un lot de médicaments, il doit
savoir \emph{pourquoi} le modèle soupçonne une falsification, puis \emph{quoi
faire}. Ce chapitre apporte ces deux réponses \textbf{pour chaque alerte}
levée par le modèle :

\begin{enumerate}
  \item l'\textbf{explication} de l'alerte par la méthode SHAP : quelles
        variables ont fait basculer le score vers « attaque » ;
  \item un \textbf{rapport de recommandations} rédigé à partir de cette
        explication : ce qui s'est passé, la cause probable, la gravité et les
        actions à mener.
\end{enumerate}

Un point de méthode conditionne tout le chapitre : le modèle est
\textbf{binaire}. Il est entraîné sur l'étiquette « normal » ou « attaque » et
ne nomme jamais l'attaque. Le type d'attaque que le rapport propose est donc
une \emph{hypothèse déduite} des variables qui ont le plus pesé dans l'alerte,
jamais une réponse du modèle. Dans l'étude, les scénarios sont connus parce que
nous les avons injectés : ils servent à \emph{mesurer} la fiabilité de cette
déduction (section~\ref{sec:expl_eval}).


% =====================================================================
\section{Explicabilité par SHAP}
\label{sec:expl_shap}

\subsection{Principe}

SHAP (\textit{SHapley Additive exPlanations}) transpose aux modèles prédictifs
la théorie des jeux coopératifs : chaque variable est un joueur, le score du
modèle est le gain à partager. La propriété essentielle est l'\textbf{additivité} :
pour une mesure $x$, le score est la valeur de base $\phi_0$ (prédiction moyenne)
augmentée des contributions $\phi_i$ de toutes les variables :

\begin{equation}
f(x) = \phi_0 + \sum_{i=1}^{n} \phi_i
\label{eq:additivite}
\end{equation}

Une valeur $\phi_i$ positive pousse la mesure vers « attaque », une valeur
négative vers « normal ». Le modèle retenu (Extra Trees, chapitre~\ref{chap4})
est un ensemble d'arbres : nous utilisons \texttt{shap.TreeExplainer}, qui
calcule ces valeurs de façon exacte en exploitant la structure des arbres.

\subsection{Importance globale et par scénario}

Pour savoir si le modèle s'appuie sur des indices cohérents, nous calculons
l'importance d'une variable $j$ comme la moyenne des valeurs absolues de ses
contributions sur $N$ mesures d'attaque non vues à l'entraînement :

\begin{equation}
I_j = \frac{1}{N}\sum_{k=1}^{N}\left|\phi_j^{(k)}\right|
\label{eq:importance}
\end{equation}

Les variables sont classées une fois sur l'ensemble des attaques (importance
globale) et une fois par scénario injecté. Ces classements sont enregistrés
dans la table \texttt{transmission\_shap} et affichés par la page
« Explicabilité SHAP » du tableau de bord.

\begin{figure}[H]
\centering
\figtrois{fig_shap_global.png}{0.85}
\caption{Importance globale des variables (moyenne des valeurs SHAP
absolues), Extra Trees.}
\label{fig:expl_shap_global}
\end{figure}

Ces classements montrent que chaque famille d'attaque est reconnue par la
variable conçue pour elle : la séquence de 2 minutes déjà transmise pour le
rejeu, l'écart-type à 1 minute pour le gel et l'injection nominale, et la forme
de la distribution sur 5 à 15 minutes (moyenne, maximum, quantiles) pour les
biais et la dérive. C'est cette correspondance entre un type d'attaque et une
signature dans les variables qui rend possible la déduction décrite plus bas.
Une précaution de lecture : l'importance SHAP n'est pas une causalité, et deux
variables corrélées se partagent l'effet.

\subsection{SHAP par alerte}
\label{sec:expl_alerte}

L'importance globale décrit le modèle, pas une alerte précise. Pour qu'un
opérateur comprenne \emph{son} alerte, nous calculons le SHAP de chaque alerte,
c'est-à-dire d'une suite continue de mesures en alarme (règle de persistance :
au moins $k$ mesures sur les $n$ dernières au-dessus du seuil). Le calcul est
le suivant :

\begin{enumerate}
  \item on prend les lignes de l'alerte ; au-delà de 6 lignes, on en retient 6
        réparties régulièrement sur sa durée, car le calcul exact sur 400 arbres
        profonds coûte environ 0,3 s par ligne ;
  \item \texttt{TreeExplainer} donne une valeur SHAP par ligne et par variable ;
  \item on fait la moyenne \emph{signée} sur les lignes retenues (et non la
        valeur absolue) : on cherche les variables qui ont poussé \emph{vers}
        « attaque » ;
  \item on garde les 6 variables dont la moyenne est la plus positive ;
  \item pour chacune, on enregistre aussi sa valeur moyenne pendant l'alerte et
        sa valeur moyenne hors attaque, prise comme référence « fonctionnement
        normal ».
\end{enumerate}

Le résultat est stocké dans la table \texttt{transmission\_alert\_shap}
(colonnes : alerte, rang, variable, valeur SHAP, valeur pendant l'alerte,
valeur normale) et renvoyé par la route \texttt{GET /api/v1/etude/timeline}
avec chaque alerte. Il est calculé à la fin d'un entraînement lancé depuis
l'interface ; pour un entraînement déjà enregistré, le script
\texttt{shap\_alertes.py} le recalcule sans réentraîner. Sur notre entraînement
(56 alertes), le calcul prend environ une minute et demie.

Le point clé est la comparaison « valeur pendant l'alerte contre valeur
normale » : elle transforme un nombre abstrait en fait lisible. Par exemple,
pour un gel du frigo, l'écart-type à 1 minute vaut 0,00 pendant l'alerte contre
3,16 en fonctionnement normal.


% =====================================================================
\section{Du SHAP au rapport de recommandations}
\label{sec:expl_rapport}

\subsection{Principe : faits mesurés, puis hypothèse, puis actions}

Le rapport d'une alerte est généré automatiquement par la fonction
\texttt{rapportAlerte} du tableau de bord. Il sépare nettement ce qui est
\emph{mesuré} de ce qui est \emph{supposé} :

\begin{itemize}
  \item la \textbf{synthèse} (mesuré) : date, durée, capteur concerné, score
        maximal et variable la plus déterminante avec sa valeur normale ;
  \item l'\textbf{hypothèse} (supposé) : la famille d'attaque probable et sa
        confiance, toujours accompagnée de la mention « à confirmer sur le
        terrain » ;
  \item la \textbf{gravité} (règle) : élevée, moyenne ou à vérifier ;
  \item les \textbf{actions} (règle métier), classées en trois niveaux :
        immédiat, vérification, suivi.
\end{itemize}

\subsection{Synthèse}

La synthèse est une phrase assemblée à partir de données lues en base :

\begin{table}[H]
\centering
\small
\renewcommand{\arraystretch}{1.3}
\caption{Origine de chaque élément de la synthèse.}
\label{tab:expl_synthese}
\begin{tabularx}{\textwidth}{@{}|>{\raggedright\arraybackslash}p{3.2cm}|>{\raggedright\arraybackslash}X|@{}}
\hline
\rowcolor{thead}
\textcolor{white}{\textbf{Élément}} & \textcolor{white}{\textbf{Origine}} \\
\hline
Date et durée & début et durée de l'alerte (table \texttt{transmission\_alerts}) \\
\rowcolor{trowa}
Capteur & capteur de l'épisode injecté ; à défaut, préfixe de la première
variable SHAP (\texttt{fridge\_}, \texttt{thermostat\_}) \\
Score maximal & plus haut score du modèle sur les lignes de l'alerte \\
\rowcolor{trowa}
Variable déterminante & première variable du SHAP de l'alerte, avec sa valeur
pendant l'alerte et sa valeur normale \\
\hline
\end{tabularx}
\end{table}

Le score de chaque mesure est la probabilité d'attaque donnée par
\texttt{predict\_proba} : pour un Extra Trees de 400 arbres, c'est la moyenne
des probabilités des arbres. Le \emph{score maximal} est le plus haut de ces
scores sur l'alerte. Il décrit le pic, pas l'ensemble de l'alerte : c'est une
limite à garder en tête (section~\ref{sec:expl_limites}).

\subsection{Famille probable et confiance}
\label{sec:expl_deduction}

Le modèle ne dit pas « ceci est un gel ». Les variables SHAP sont donc
regroupées en trois familles selon leur nature :

\begin{table}[H]
\centering
\small
\renewcommand{\arraystretch}{1.3}
\caption{Familles d'attaque et variables associées.}
\label{tab:expl_familles}
\begin{tabularx}{\textwidth}{@{}|>{\raggedright\arraybackslash}p{3.2cm}|>{\raggedright\arraybackslash}X|>{\raggedright\arraybackslash}X|@{}}
\hline
\rowcolor{thead}
\textcolor{white}{\textbf{Famille}} & \textcolor{white}{\textbf{Variables}} &
\textcolor{white}{\textbf{Scénarios simulés}} \\
\hline
Signal figé & écart-type, variation absolue, variation d'une mesure à l'autre
& gel du frigo, gel du thermostat, injection nominale \\
\rowcolor{trowa}
Séquence rejouée & séquence de 2 min déjà transmise & rejeu du frigo \\
Décalage de niveau & moyenne, minimum, maximum, quantiles, écart entre
capteurs & biais ($-0{,}5$, $-1$, $-2$\,\textdegree C, thermostat), dérive \\
\hline
\end{tabularx}
\end{table}

La famille est ensuite choisie par trois règles, appliquées dans cet ordre :

\begin{enumerate}
  \item \textbf{Rejeu.} Si la variable « séquence de 2 min déjà transmise » est
        active pendant l'alerte (elle vaut 0 en fonctionnement normal) et pousse
        vers « attaque », la famille est « séquence rejouée ». La confiance
        affichée est l'écart à la valeur normale.
  \item \textbf{Signal figé.} Si un écart-type ou une variation absolue tombe à
        moins de 25\,\% de sa valeur normale, la famille est « signal figé ». La
        confiance est l'écart relatif à la valeur normale (0,00 contre 3,16
        donne 100\,\%).
  \item \textbf{Vote pondéré.} Sinon, on additionne l'effet SHAP positif de
        chaque variable par famille. La famille dominante est retenue si elle
        pèse au moins 50\,\% du total, avec cette part pour confiance ; sinon la
        réponse est « indéterminé » : le rapport ne force pas une réponse.
\end{enumerate}

Les deux premières règles ont été ajoutées après un constat : un gel fige la
mesure mais fait \emph{aussi} baisser la moyenne, et un rejeu copie des valeurs
normales passées. Dans les deux cas, le simple vote pondéré donnait « décalage
de niveau » ou « indéterminé » à tort, parce que des variables de niveau
cumulaient plus d'effet que l'écart-type ou la séquence répétée. Le seuil de
25\,\% a été choisi sur un plateau : tout seuil entre 20 et 30\,\% donne le même
résultat sur nos données. La confiance affichée est un indicateur de netteté du
signal, pas une probabilité calibrée.

\subsection{Gravité}

La gravité aide l'opérateur à ordonner ses interventions. Elle n'utilise que
deux informations que l'opérateur connaît au moment de l'alerte : sa
\textbf{durée} et son \textbf{score maximal}, c'est-à-dire le plus haut score de
probabilité d'attaque donné par le modèle sur les mesures de l'alerte. Elle
n'utilise jamais l'étiquette « vraie » ou « fausse », qui n'est connue dans
l'étude que parce que les attaques sont injectées. La règle est la suivante :

\begin{table}[H]
\centering
\small
\renewcommand{\arraystretch}{1.3}
\caption{Règle de gravité d'une alerte.}
\label{tab:expl_gravite}
\begin{tabular}{|l|l|}
\hline
\rowcolor{thead}
\textcolor{white}{\textbf{Gravité}} & \textcolor{white}{\textbf{Condition}} \\
\hline
Élevée & durée d'au moins 15 min, ou score maximal d'au moins 0,95 \\
\rowcolor{trowa}
À vérifier & durée de moins de 5 min et score maximal inférieur à 0,80 \\
Moyenne & tous les autres cas \\
\hline
\end{tabular}
\end{table}

L'idée est qu'une manipulation longue expose plus longtemps les produits, qu'un
score très élevé traduit un modèle très sûr de lui, et qu'une alerte courte au
score peu élevé est plus souvent un faux déclenchement. Sur les 56 alertes de
l'entraînement de référence, cette règle classe 49 des 51 alertes vraies en
gravité élevée et 2 en moyenne ; aucune alerte vraie n'est « à vérifier », contre
3 des 5 fausses alertes. L'échantillon de fausses alertes est très petit. Les
seuils de 15 min, 0,95, 5 min et 0,80 sont des \textbf{choix de conception} : ils
ne résultent ni d'une étude ni d'une norme pharmaceutique et doivent être
ajustés avec un responsable qualité. La gravité ne tient compte ni du type de produit, ni de l'écart réel à
la plage de conservation (2--8\,\textdegree C).

\subsection{Actions suggérées}

Les actions sont classées dans l'ordre d'un incident qualité : protéger,
comprendre, tracer.

\begin{table}[H]
\centering
\small
\renewcommand{\arraystretch}{1.35}
\caption{Actions suggérées dans le rapport d'une alerte vraie.}
\label{tab:expl_actions}
\begin{tabularx}{\textwidth}{@{}|>{\raggedright\arraybackslash}p{2.6cm}|>{\raggedright\arraybackslash}X|@{}}
\hline
\rowcolor{thead}
\textcolor{white}{\textbf{Niveau}} & \textcolor{white}{\textbf{Contenu}} \\
\hline
Immédiat & mettre les produits concernés en quarantaine ; contrôler la
température avec un thermomètre de référence, sans se fier au capteur concerné
\\
\hline
\rowcolor{trowa}
Vérification (signal figé) & vérifier que le capteur et la liaison ne sont pas
bloqués ; comparer avec le second capteur \\
\hline
Vérification (rejeu) & vérifier l'intégrité de la liaison
capteur--passerelle ; contrôler l'horodatage des mesures reçues \\
\hline
\rowcolor{trowa}
Vérification (niveau) & vérifier l'étalonnage et le groupe froid ; chercher
une cause réelle (porte ouverte, panne de froid) avant de conclure à une
falsification \\
\hline
Suivi & consigner l'incident et prévenir le responsable qualité \\
\hline
\end{tabularx}
\end{table}

Pour une alerte d'un type « indéterminé », seules les actions immédiates et de
suivi sont proposées. Pour une \textbf{fausse alerte connue} (cas qui n'existe
que dans la page d'évaluation « Détection vs réalité », où l'on compare le
modèle à la réalité), le rapport ne demande pas de quarantaine : il propose de
contrôler la température, puis de consigner la fausse alerte pour ajuster le
seuil. Un opérateur en exploitation ne connaît pas cette étiquette : il s'appuie
sur la gravité « à vérifier ». Ces actions relèvent du bon sens
métier de la chaîne du froid ; elles sont présentées comme des suggestions que
l'opérateur valide.

\subsection{Affichage dans le tableau de bord}

Le rapport est intégré au tableau « score du modèle et attaques réelles » de la
page « Détection vs réalité ». Ce tableau compte une ligne par minute ; une
alerte de 22 minutes couvrirait donc 22 lignes. Pour ne pas répéter la même
recommandation, elle n'apparaît que sur la \textbf{première ligne de chaque
alerte}, avec sa gravité, sa famille probable et un bouton « Rapport ». Les
lignes suivantes de la même alerte renvoient à la première. Le bouton déroule
la synthèse, l'hypothèse, les variables SHAP et les actions.

\begin{figure}[H]
\centering
\figtrois{capture_detection_rapport.png}{0.95}
\caption{Rapport de recommandation d'une alerte dans la page « Détection vs
réalité ».}
\label{fig:expl_capture_rapport}
\end{figure}

\subsection{Deux exemples}

\textbf{Gel du frigo} (alerte de 22,2 min, score maximal 1,00). Variables
dominantes : écart-type à 1 min (0,00 contre 3,16 en fonctionnement normal),
moyenne à 5 min (3,79 contre 7,13), maximum à 5 min (4,89 contre 13,38). La
règle « signal figé » s'applique (confiance 100\,\%), la gravité est élevée, et
la vérification demandée porte sur un capteur ou une liaison bloqués.

\textbf{Rejeu du frigo} (alerte de 15,6 min, score maximal 0,96). Variable
dominante : séquence de 2 min déjà transmise (1,00 contre 0,00). La règle
« rejeu » s'applique (confiance 100\,\%), la gravité est élevée, et la
vérification demandée porte sur la liaison capteur--passerelle et les
horodatages.


% =====================================================================
\section{Évaluation de la famille déduite}
\label{sec:expl_eval}

Puisque les scénarios sont injectés, nous pouvons mesurer si la famille déduite
du SHAP correspond à la famille réellement simulée. L'évaluation porte sur les
alertes vraies de la répétition 0 de l'entraînement n$^\circ$\,2 (Extra Trees,
20 répétitions) : 56 alertes au total, dont 51 vraies avec un scénario connu.

\begin{table}[H]
\centering
\small
\renewcommand{\arraystretch}{1.3}
\caption{Famille déduite correcte, selon les règles employées (51 alertes
vraies).}
\label{tab:expl_eval}
\begin{tabular}{|l|c|}
\hline
\rowcolor{thead}
\textcolor{white}{\textbf{Méthode}} & \textcolor{white}{\textbf{Alertes correctes}} \\
\hline
Vote pondéré seul & 44 sur 51 \\
\rowcolor{trowa}
Vote + règle du signal figé & 49 sur 51 \\
Vote + règles du signal figé et du rejeu & 51 sur 51 \\
\hline
\end{tabular}
\end{table}

Ce résultat doit être lu avec prudence, pour trois raisons.

\begin{enumerate}
  \item \textbf{Peu d'alertes hors apprentissage.} Sur les 51 alertes, 43
        tombent sur des heures d'entraînement et seulement 8 sur des heures de
        validation (4) et de test (4). Sur ces 8 alertes, le vote seul en
        classe 7 correctement et les trois règles les 8. La
        proportion est encourageante, mais l'échantillon est petit.
  \item \textbf{Règles réglées sur ces mêmes alertes.} Le seuil de 25\,\% et la
        règle du rejeu ont été construits en observant ces 51 alertes : le
        résultat de 51 sur 51 mesure l'adéquation des règles à ce run, pas leur
        capacité à généraliser. Une validation honnête demanderait un autre
        entraînement, avec d'autres épisodes injectés.
  \item \textbf{Attaques simulées.} Les attaques sont injectées dans des traces
        réelles TON\_IoT avec des signatures connues. Une attaque réelle
        pourrait avoir une autre signature.
\end{enumerate}

Nous formulons donc ainsi le résultat : « les règles ont été construites et
vérifiées sur les 51 alertes de l'entraînement de référence ; la généralisation
reste à confirmer sur de nouveaux épisodes ».


% =====================================================================
\section{Limites}
\label{sec:expl_limites}

\begin{itemize}
  \item \textbf{Hypothèse, pas diagnostic.} Le modèle est binaire ; la famille
        n'est qu'une déduction, à confirmer par une vérification sur place.
  \item \textbf{Échantillonnage du SHAP.} Une alerte longue est expliquée par 6
        lignes réparties sur sa durée, pour limiter le temps de calcul.
  \item \textbf{Score maximal.} Il décrit le pic de l'alerte, pas sa tenue
        dans la durée ; le score moyen serait un complément utile.
  \item \textbf{Seuils de gravité.} 15 min, 0,95, 5 min et 0,80 sont des choix de
        conception, sans lien avec le type de produit ni avec l'écart réel à la
        plage de conservation. Ils ont été vérifiés sur 56 alertes seulement,
        dont 5 fausses.
  \item \textbf{Familles proches.} Un biais et une dérive déplacent les mêmes
        variables : les indices ne permettent pas de les distinguer, ils sont
        regroupés dans la même famille « décalage de niveau ».
  \item \textbf{Temps réel simulé.} Les alertes de ce chapitre viennent de
        traces rejouées, pas d'un flux de capteurs réels.
\end{itemize}

Les règles de conduite valables pour toute alerte restent : vérifier d'abord
sur place avec un thermomètre de contrôle, ne jamais libérer un lot sur la seule
base d'une mesure contestée, et journaliser chaque alerte, la vérification
faite et la décision prise.


% =====================================================================
\section{Conclusion}

Ce chapitre a relié le modèle de détection à l'action de l'opérateur. Pour
chaque alerte, l'application calcule un SHAP propre à cette alerte, en déduit
une famille probable de manipulation par trois règles simples, puis rédige un
rapport : une synthèse de faits mesurés, une hypothèse avec sa confiance, une
gravité et des actions classées par priorité. Le rapport n'apparaît qu'une
fois par alerte, dans le tableau de détection, et sépare toujours ce qui est
mesuré de ce qui est supposé.

Sur les 51 alertes vraies de l'entraînement de référence, la famille déduite
correspond au scénario injecté dans tous les cas après ajout des deux règles
prioritaires, contre 44 sur 51 avec un vote simple. Ce chiffre reste mesuré sur
les données qui ont servi à régler les règles, et seules 8 alertes sont hors
apprentissage : la généralisation reste à confirmer. La perspective naturelle
est un modèle multiclasse, entraîné à prédire directement la famille, au prix
d'un nouvel entraînement et d'une nouvelle évaluation. Le chapitre suivant
décrit la réalisation de l'API et du tableau de bord qui exposent ces
résultats.
